\uuid{2D0M}
\titre{Nombres premiers}
\auteur{bodin ; croizat ; sacre}
\organisation{exo7}
\licence{CC BY-NC-SA 4.0}
\datecreate{2022-03-08}
\niveau{L1}
\module{Algèbre}
\chapitre{Arithmétique dans Z}
\sousChapitre{Nombres premiers, nombres premiers entre eux}
\difficulte{2}
\theme{Nombres premiers}
\exo7qcmid{hilisit.arithmetique.035}

\contenu{
\begin{question}
\qclassification{103.04}
 Soit $p$ un nombre premier tel que $10 < p < 100$. On note $A$ le chiffre des dizaines et $B$ le chiffre des unités de l'écriture décimale de $p$. Quelles sont les affirmations vraies ?
    \begin{answers}
        \good{$A$ peut être pair.}
        \bad{$B$ peut être pair.}
        \good{On peut avoir $A = B$.}
        \bad{On peut avoir $B = 9 - A$.}
    \end{answers}
    \begin{explanations}
    Pour $p = 23$ qui est premier, on a bien $A = 2$ qui est pair.\\
    Si le chiffre des unités $B$ est pair, alors $p$ est pair ce qui est impossible pour un nombre premier $\geq 3$.\\
    Pour $p=11$ premier, on a bien $A=B$ (les autres nombres avec deux chiffres identiques sont justement les multiples de $11$, et ne sont donc pas premiers).\\
    Si $B = 9-A$, alors la somme des chiffres de $p$ vaut $A + B = 9$ : ainsi $p$ est divisible par $9$, ce qui contredit sa primalité.
    \end{explanations}
\end{question}
}
